\section{Infix to Postfix Conversion}
\label{sec:sq-infix-to-postfix}

\noindent We now turn to a group of six closely related problems:
converting between the three ways of writing an arithmetic expression.
\textbf{Infix} places operators between operands ($a+b$) and needs
parentheses and precedence rules to remove ambiguity. \textbf{Postfix}
(operators after their operands, $ab+$) and \textbf{prefix} (operators
before, $+ab$) are both unambiguous without any parentheses, which is
exactly why calculators and compilers prefer them internally -- and why a
stack-based scan can evaluate or convert them in a single pass.

\problemstatement{Given an infix expression (operands, the binary
operators \texttt{+ - * / \textasciicircum}, and parentheses), convert it
to postfix notation.}

\begin{patternbox}
\emph{``Convert an expression respecting operator precedence and
parentheses''} is the Shunting-Yard algorithm's signature: a stack holds
operators \emph{waiting} to be placed into the output, and the central
decision at every operator is \emph{``do any waiting operators have
\textbf{equal or higher} precedence than the one I just read, and if so,
should they come out first?''} -- exactly the kind of ordering decision a
stack resolves naturally.
\end{patternbox}

\subsection*{Understanding the problem carefully}

In postfix, an operator appears \emph{immediately after} both of its
operands are fully written -- which means an operator can only be emitted
to the output once we are certain nothing of higher (or equal, for
left-associative operators) precedence is still waiting to bind tighter
around it. Parentheses simply force this resolution explicitly: everything
inside a $(\ldots)$ must be fully resolved into postfix before the
parenthesis is allowed to close.

\begin{ideabox}[Precedence and Associativity]
Assign precedence: \texttt{\textasciicircum} (highest, right-associative)
$>$ \texttt{* /} (left-associative) $>$ \texttt{+ -} (lowest,
left-associative). When deciding whether to pop an operator already on
the stack before pushing a new one: pop if the stacked operator has
\textbf{strictly higher} precedence, or \textbf{equal} precedence \emph{and}
the new operator is left-associative (for a right-associative operator
like \texttt{\textasciicircum}, equal precedence does \emph{not} trigger a
pop, since \texttt{\textasciicircum} groups right-to-left:
$2\texttt{\textasciicircum}3\texttt{\textasciicircum}2$ means
$2\texttt{\textasciicircum}(3\texttt{\textasciicircum}2)$, not
$(2\texttt{\textasciicircum}3)\texttt{\textasciicircum}2$).
\end{ideabox}

\subsection*{Why popping by precedence produces correct postfix}

An operator should only be written to postfix output once everything that
needs to bind \emph{more tightly} around its operands has already been
resolved. If the stack's top has strictly higher precedence, it must
finish binding its own operands first (it is ``more urgent''), so we pop
it to the output before pushing the new, lower-precedence operator. If
precedences are equal and the operator is left-associative, the earlier
(already-stacked) operator should also be resolved first, since
left-associativity means we group left to right -- e.g.\ $a-b-c$ means
$(a-b)-c$, so the first \texttt{-} must close over $a,b$ before the second
\texttt{-} is considered. Right-associative operators reverse this rule
specifically to preserve right-to-left grouping.

\subsection*{Algorithm}

\begin{enumerate}
  \item Scan left to right. For each token:
  \begin{enumerate}
    \item Operand: append to output.
    \item \texttt{'('}: push.
    \item \texttt{')'}: pop and append to output until \texttt{'('} is
          popped (discard both parentheses).
    \item Operator $op$: while the stack's top is an operator with
          higher precedence, or equal precedence with $op$
          left-associative: pop and append to output. Then push $op$.
  \end{enumerate}
  \item After the scan, pop all remaining operators to the output.
\end{enumerate}

\begin{algorithm}[H]
\caption{Infix to Postfix}
\begin{algorithmic}[1]
\Procedure{InfixToPostfix}{$s$}
  \State $\textit{output} \gets$ empty, $\textit{stack} \gets$ empty
  \For{each token $t$ in $s$}
    \If{$t$ is an operand} append $t$ to \textit{output}
    \ElsIf{$t = ($} push $t$
    \ElsIf{$t = )$}
      \While{top $\ne$ (} pop, append to \textit{output} \EndWhile
      \State pop the \texttt{(}
    \Else
      \While{\textit{stack} not empty \textbf{and} top has higher precedence, \textbf{or} (equal precedence \textbf{and} $t$ is left-assoc)}
        \State pop, append to \textit{output}
      \EndWhile
      \State push $t$
    \EndIf
  \EndFor
  \While{\textit{stack} not empty} pop, append to \textit{output} \EndWhile
  \State \Return \textit{output}
\EndProcedure
\end{algorithmic}
\end{algorithm}

\subsection*{Dry run}

\dryrun{Take $s = \texttt{"a+b*c"}$.}

\begin{itemize}
  \item \texttt{a}: output $=\texttt{"a"}$.
  \item \texttt{+}: stack empty, push. Stack $=[+]$.
  \item \texttt{b}: output $=\texttt{"ab"}$.
  \item \texttt{*}: top is \texttt{+} (lower precedence than \texttt{*}),
        no pop. Push. Stack $=[+,*]$.
  \item \texttt{c}: output $=\texttt{"abc"}$.
  \item End of scan: pop all. Pop \texttt{*}: output
        $=\texttt{"abc*"}$. Pop \texttt{+}: output $=\texttt{"abc*+"}$.
\end{itemize}

Result: \texttt{"abc*+"}, correctly reflecting that \texttt{*} binds
tighter than \texttt{+} (i.e.\ $a+(b*c)$).

\subsection*{C++ implementation}

\begin{lstlisting}
#include <bits/stdc++.h>
using namespace std;

int precedence(char op) {
    if (op == '^') return 3;
    if (op == '*' || op == '/') return 2;
    if (op == '+' || op == '-') return 1;
    return 0;
}

bool isRightAssoc(char op) { return op == '^'; }

string infixToPostfix(string s) {
    string output;
    stack<char> st;

    for (char c : s) {
        if (isalnum(c)) {
            output += c;
        } else if (c == '(') {
            st.push(c);
        } else if (c == ')') {
            while (st.top() != '(') {
                output += st.top();
                st.pop();
            }
            st.pop(); // remove '('
        } else { // operator
            while (!st.empty() && st.top() != '(' &&
                   (precedence(st.top()) > precedence(c) ||
                    (precedence(st.top()) == precedence(c) && !isRightAssoc(c)))) {
                output += st.top();
                st.pop();
            }
            st.push(c);
        }
    }
    while (!st.empty()) {
        output += st.top();
        st.pop();
    }
    return output;
}

int main() {
    cout << infixToPostfix("a+b*c") << endl;
    return 0;
}
\end{lstlisting}

\begin{complexitybox}
\textbf{Time Complexity.} Each token is pushed and popped at most once,
giving
\[
  O(n).
\]
\textbf{Space Complexity.} $O(n)$ for the stack and output in the worst
case.
\end{complexitybox}

\begin{takeawaybox}
The precedence-and-associativity comparison rule derived here is reused,
completely unchanged, in Infix to Prefix
(Section~\ref{sec:sq-infix-to-prefix}) -- only the direction of scanning
and the meaning of the output change.
\end{takeawaybox}
